\section{Тест производительности}

Тест производительности представляет из себя сравнение адаптированного
алгоритма Ахо~--~Корасик с наивным алгоритмом поиска (полным перебором
позиций). Сравнение проводится на двух контрастных сценариях,
демонстрирующих сильные и слабые стороны реализованного алгоритма.

\textbf{Сценарий 1 (Ахо~--~Корасик выигрывает).}
Образец содержит несколько джокеров в начале и длинный твёрдый суффикс:
\texttt{? ? ? ? 4 2 3 9 1}. Текст состоит из случайных чисел от $0$ до
$99$; твёрдый суффикс искусственно вставляется в текст достаточно
редко (примерно раз в $10\,000$ символов). В таких условиях наивный
алгоритм вынужден на каждой позиции проверять до $9$ символов образца,
в то время как автомат Ахо~--~Корасик быстро проскакивает несовпадающие
участки и проверяет только те позиции, где встретился суффикс.

\textbf{Сценарий 2 (Ахо~--~Корасик проигрывает).}
Образец состоит из часто повторяющихся коротких фрагментов,
разделённых джокерами: \texttt{1 ? 1 ? 1}. Текст на $90\%$ состоит из
единиц и на $10\%$ из двоек. Каждая единица в тексте порождает
вхождение твёрдого фрагмента \texttt{1}, что приводит к лавинному
росту количества сохранённых позиций ($Z \approx 0{,}9 \cdot n \cdot 3$).
Наивный алгоритм при этом не несёт накладных расходов на хранение
и проверку позиций-кандидатов и работает быстрее.

Для каждого сценария замерялось чистое время поиска (без учёта
считывания входных данных и вывода результатов). Тестирование
проводилось на текстах размера $10\,000$, $100\,000$ и $1\,000\,000$
чисел.

\begin{alltt}
root@6ec7ac378e08:/workspaces/discran-labs/laba-4# ./banch.out
pattern: ? ? ? ? 4 2 3 9 1
n           Naive (ms)    Aho (ms)      Speedup
--------------------------------------------------
10000       0.09          0.06          1.61
100000      1.12          0.70          1.60
1000000     10.48         5.78          1.81

pattern: 1 ? 1 ? 1
n           Naive (ms)    Aho (ms)      Speedup
--------------------------------------------------
10000       0.12          0.26          0.45
100000      1.35          2.85          0.47
1000000     35.95         80.56         0.45
\end{alltt}

Как видно, в первом сценарии адаптированный алгоритм Ахо~--~Корасик
стабильно обходит наивный перебор примерно в $1{,}6$--$1{,}8$ раза.
Относительный выигрыш не зависит от размера текста, что согласуется
с теоретической оценкой: при редких вхождениях фрагментов сложность
близка к линейной $O(n + m)$, тогда как наивный алгоритм всегда
выполняет $O(n \cdot m)$ сравнений.

Во втором сценарии автомат Ахо~--~Корасик проигрывает наивному
алгоритму примерно в $2{,}2$ раза. Проигрыш вызван накладными
расходами на сохранение и последующую обработку большого числа
вхождений коротких фрагментов ($Z \approx 0{,}9 \cdot n \cdot 3$).
Наивный алгоритм не выделяет дополнительной памяти под
промежуточные результаты и в данном случае оказывается
эффективнее.

Таким образом, адаптированный алгоритм Ахо~--~Корасик является
предпочтительным для образцов с длинными твёрдыми фрагментами,
которые редко встречаются в тексте. В случае же, когда образец
состоит из коротких часто встречающихся фрагментов, наивный
алгоритм может оказаться более производительным из-за отсутствия
накладных расходов на хранение промежуточных данных.
\pagebreak